\documentclass{beamer}

\usepackage[utf8]{inputenc}  % Para asegurar que usa UTF-8
\usepackage[T1]{fontenc}     % Para mejorar el soporte de caracteres acentuados

\usepackage[spanish]{babel}
\usepackage{listings}      % Para código fuente
\usepackage{xcolor}        % Colores para resaltar código
\usepackage{hyperref}      % Enlaces

\usepackage{lmodern}
\usepackage{textcomp}

%--- Tema y colores de beamer ---
\usetheme{Madrid}

% Definimos colores personalizados (opcional)
\definecolor{myBlue}{RGB}{25,25,112}  % MidnightBlue
\definecolor{myGray}{RGB}{240,240,240}

% Configuramos un fondo para la diapositiva de cambio de sección (opcional)
% Ejemplo: color de fondo gris claro
\setbeamercolor{section page}{bg=myGray, fg=myBlue}

% Podemos usar \AtBeginSection para mostrar una diapositiva indicadora
\AtBeginSection[]{
  \begin{frame}[plain]  % "plain" para que no aparezca encabezado/footers
    \centering
    % Podemos usar un beamercolorbox personalizado:
    \begin{beamercolorbox}[sep=10pt, center, shadow=true, rounded=true]{section page}
      \usebeamerfont{title}\Huge \textbf{\insertsectionhead}
    \end{beamercolorbox}
  \end{frame}
}

% Opcional: definimos un estilo para el código
\lstset{
    inputencoding=utf8,     % Asegurar codificación UTF-8
  backgroundcolor=\color{white},
  basicstyle=\small\ttfamily,
  breaklines=true,
  keywordstyle=\bfseries\color{blue},
  commentstyle=\itshape\color{gray},
  stringstyle=\color{red},
  showstringspaces=false,
    extendedchars=true,     % Permitir caracteres extendidos
    frame=single,           % Marco alrededor del código
    language=Java,           % Especificar el lenguaje si aplica
        literate={á}{{\'a}}1 {é}{{\'e}}1 {í}{{\'i}}1 {ó}{{\'o}}1 {ú}{{\'u}}1
             {Á}{{\'A}}1 {É}{{\'E}}1 {Í}{{\'I}}1 {Ó}{{\'O}}1 {Ú}{{\'U}}1
             {ñ}{{\~n}}1 {Ñ}{{\~N}}1
}

\begin{document}

%-------------------------------------
\begin{frame}{Índice (1/3)}
  \tableofcontents[sections={1-3}]
\end{frame}

\begin{frame}{Índice (2/3)}
  \tableofcontents[sections={4-5}]
\end{frame}


\begin{frame}{Índice (3/3)}
  \tableofcontents[sections={6-8}]
\end{frame}


%=====================================
\section{Introducción a Java}

%-------------------------------------
\subsection{Historia y filosofía de Java}

\begin{frame}{1.1.1. Orígenes de Java}
\textbf{Creador principal}: James Gosling (Sun Microsystems) \\
\textbf{Año de concepción}: Alrededor de 1991 \\
\textbf{Nombre original}: “Oak” \\
\textbf{¿Por qué “Java”?} \\
\begin{itemize}
  \item Se dice que el nombre proviene de una marca de café popular entre los desarrolladores.
  \item El ícono de la taza de café humeante alude a esta anécdota.
\end{itemize}

\textbf{Anécdota}:  
Java inicialmente se pensó para programar dispositivos electrónicos (TV, mandos). Con el auge de Internet a mediados de los 90, pasó a aplicaciones web/escritorio y se hizo muy popular.
\end{frame}

%-------------------------------------
\begin{frame}{1.1.2. Filosofía de Java}
\begin{itemize}
  \item \textbf{Simplicidad}: Sintaxis familiar a C/C++.
  \item \textbf{Portabilidad} (WORA): “Write Once, Run Anywhere”.
  \item \textbf{Robustez}: Manejo automático de memoria (GC), detección temprana de errores.
  \item \textbf{Seguridad}: “Sandbox” que limita acciones peligrosas.
  \item \textbf{Orientación a objetos}: Clases y objetos como base del lenguaje.
\end{itemize}

\textbf{Reflexión}:  
¿Por qué la portabilidad WORA fue tan revolucionaria a mediados de los 90?  
\end{frame}

%-------------------------------------
\subsection{Características de Java}

\begin{frame}{1.2. Características de Java}
\begin{enumerate}
  \item \textbf{Orientación a objetos}.
  \item \textbf{Manejo automático de memoria} (Garbage Collector).
  \item \textbf{Elevada portabilidad}: Bytecode interpretado por la JVM.
  \item \textbf{Seguridad}: Modelo estricto para ejecución controlada.
  \item \textbf{Extensa biblioteca estándar} (Java API).
  \item \textbf{Multihilo}: Soporte integrado (el año que viene).
  \item \textbf{Relativa simplicidad} frente a C/C++: sin punteros aritméticos, etc.
\end{enumerate}

\textbf{Ejemplo comparativo (C vs Java)}:  
En C, `malloc` / `free` puede generar fugas de memoria. En Java, `new` + GC simplifica la gestión.
\end{frame}

%-------------------------------------
\subsection{Diferencias clave entre C y Java}

\begin{frame}{1.3. Diferencias clave con C}
\begin{itemize}
  \item \textbf{No hay punteros “explícitos”}: Sólo referencias.
  \item \textbf{Estructuras vs Clases}: En C, `struct`; en Java, todo gira en torno a clases.
  \item \textbf{Sintaxis similar, pero no igual}: Booleanos estrictos (`true`/`false`).
  \item \textbf{No hay `main()` global}: Todo debe estar en una clase con `main(String[] args)`.
  \item \textbf{Manejo de memoria}: `malloc`/`free` en C vs `new`/GC en Java.
  \item \textbf{Biblioteca estándar} muy extensa en Java.
  \item \textbf{Compilación a bytecode}: En C → binario, en Java → `.class` interpretado por la JVM.
\end{itemize}

\textbf{Anécdota}:  
La falta de punteros explícitos frustra a programadores de C, pero mejora la seguridad.
\end{frame}

%-------------------------------------
\subsection{La JVM y el bytecode}

\begin{frame}{1.4. El concepto de la JVM}
\textbf{¿Qué es la JVM?}
\begin{itemize}
  \item Interpreta/compila en caliente (JIT) el bytecode.
  \item Cada SO o plataforma tiene su propia JVM.
\end{itemize}

\textbf{Bytecode como intermedio}
\begin{itemize}
  \item `javac` compila `.java` a `.class`.
  \item Independiente de arquitectura. La JVM lo ejecuta.
\end{itemize}

\textbf{Ventajas}
\begin{itemize}
  \item Portabilidad total.
  \item Seguridad (sandbox).
  \item Optimización con JIT.
  \item Facilidad para depuración.
\end{itemize}
\end{frame}

%-------------------------------------
\subsection{JDK y JRE}

\begin{frame}{1.5. JDK (Java Development Kit) y JRE (Java Runtime Environment)}
\textbf{JRE}:
\begin{itemize}
  \item Contiene la JVM y bibliotecas para \textit{ejecutar} aplicaciones.
  \item No incluye `javac` (sin herramientas de desarrollo).
\end{itemize}

\textbf{JDK}:
\begin{itemize}
  \item Incluye JRE + compilador `javac`, `jar`, `javadoc`, etc.
  \item Necesario para \textit{programar} en Java.
\end{itemize}

\textbf{Sugerencia de instalación}:  
Actualmente se suele usar directamente el JDK (e.g. OpenJDK, Oracle JDK, etc.).
\end{frame}

%=====================================
\section{Entorno de desarrollo}

%-------------------------------------
\subsection{Instalación de Java (JDK)}

\begin{frame}{2.1.1. ¿Qué necesitamos instalar?}
\textbf{JDK}:  
\begin{itemize}
  \item Incluye compilador, JVM y herramientas.
\end{itemize}

\textbf{Versiones populares}:
\begin{itemize}
  \item OpenJDK (libre y abierto).
  \item Oracle JDK (licencia propietaria).
  \item Distribuciones de terceros (AdoptOpenJDK, Corretto, Zulu...).
\end{itemize}

\textbf{Ejercicio}:  
Verifica la versión instalada con \texttt{java -version} y discute las diferencias.
\end{frame}

\begin{frame}{2.1.2. Dónde conseguirlo}
\begin{itemize}
  \item \href{https://adoptium.net/}{Adoptium} (OpenJDK).
  \item \href{https://www.oracle.com/java/technologies/downloads/}{Oracle} (Oracle JDK).
\end{itemize}

\textbf{Versiones LTS}: Java 11, Java 17, etc.  
Se recomiendan para proyectos a largo plazo.

\textbf{Tip práctico}:  
Muestra cómo instalar en Windows/Mac/Linux, y luego compilar un “Hola Mundo”.
\end{frame}

%-------------------------------------
\subsection{Configuración de variables de entorno}

\begin{frame}{2.2. PATH y JAVA\_HOME}
\textbf{PATH}: Permite usar \texttt{javac} y \texttt{java} en la terminal.\\
\textbf{JAVA\_HOME}: Almacena la ruta del JDK; algunas herramientas lo necesitan.

\textbf{Windows}:
\begin{enumerate}
  \item Panel de Control → Variables de entorno.
  \item Crea \texttt{JAVA\_HOME} apuntando al directorio del JDK.
  \item Agrega \texttt{\%JAVA\_HOME\%\\bin} a \texttt{PATH}.
\end{enumerate}

\textbf{macOS / Linux}:
\begin{enumerate}
  \item Editar \texttt{\$HOME/.bashrc} o similar.
  \item \texttt{export JAVA\_HOME=/ruta/al/jdk}
  \item \texttt{export PATH=\$JAVA\_HOME/bin:\$PATH}
\end{enumerate}
\end{frame}

%-------------------------------------
\subsection{Herramientas de desarrollo}

\begin{frame}{2.3. \texttt{javac}, \texttt{java}, \texttt{jar}}
\begin{itemize}
  \item \textbf{javac}: Compilador. Convierte \texttt{.java} en \texttt{.class}.
  \item \textbf{java}: Intérprete. Ejecuta el \texttt{.class} en la JVM.
  \item \textbf{jar}: Empaqueta en archivos \texttt{.jar}.
\end{itemize}

\textbf{Uso básico}:
\begin{itemize}
  \item \texttt{javac MiPrograma.java}
  \item \texttt{java MiPrograma}
  \item \texttt{jar cf MiPrograma.jar MiPrograma.class}
\end{itemize}

\textbf{Anécdota}:  
El clásico “Hola Mundo” se compila y ejecuta en dos pasos.  
\end{frame}

%-------------------------------------
\subsection{Uso de IDE vs editores de texto}

\begin{frame}{2.4. IDEs: Eclipse, IntelliJ, NetBeans...}
\textbf{Ventajas de un IDE}:
\begin{itemize}
  \item Autocompletado, resaltado de errores, depuración integrada.
  \item Manejo de proyectos grandes (módulos, bibliotecas).
\end{itemize}

\textbf{Desventajas}:
\begin{itemize}
  \item Ocultan detalles de compilación.
  \item Requiere más recursos del sistema.
\end{itemize}

\textbf{Edición manual + terminal}:
\begin{itemize}
  \item Bueno para aprender el flujo de compilación.
  \item Alternativas ligeras: VSCode, Vim + plugins, etc.
\end{itemize}
\end{frame}

%-------------------------------------
\subsection{Primer programa “Hola Mundo”}

\begin{frame}[fragile]{2.5. Hola Mundo}
\textbf{Ejemplo}:
\begin{lstlisting}[language=Java]
public class HolaMundo {
    public static void main(String[] args) {
        System.out.println("Hola, mundo!");
    }
}
\end{lstlisting}

\textbf{Claves}:
\begin{itemize}
  \item \texttt{public class} coincide con el nombre del archivo (\texttt{HolaMundo.java}).
  \item \texttt{main(String[] args)} es el punto de entrada.
  \item \texttt{System.out.println} imprime en consola.
\end{itemize}

\textbf{Pasos}:
\begin{enumerate}
  \item \texttt{javac HolaMundo.java}
  \item \texttt{java HolaMundo}
\end{enumerate}
\end{frame}

%=====================================
\section{Fundamentos de la sintaxis de Java}

%-------------------------------------
\subsection{Estructura básica de un programa}

\begin{frame}[fragile]{3.1. Estructura}
\begin{lstlisting}[language=Java]
public class MiClase {
    // Atributos
    // Métodos
    public static void main(String[] args) {
        System.out.println("Hola, Java");
    }
}
\end{lstlisting}

\textbf{En Java}:
\begin{itemize}
  \item Todo dentro de una clase.
  \item El archivo \texttt{MiClase.java} contiene una \texttt{public class MiClase}.
  \item No hay funciones “sueltas” como en C.
\end{itemize}
\end{frame}

%-------------------------------------
\subsection{Nombrado de clases y métodos}

\begin{frame}{3.2. Convenciones de nombres}
\begin{itemize}
  \item \textbf{Clases}: PascalCase (p.ej. \texttt{MiClaseBonita}).
  \item \textbf{Métodos y variables}: camelCase (p.ej. \texttt{sumarNumeros}).
  \item \textbf{Constantes}: MAYÚSCULAS con guiones bajos (p.ej. \texttt{MAX\_VALUE}).
\end{itemize}

\textbf{Ejemplo}:
\begin{itemize}
  \item \texttt{public class Calculadora \{...\}}
  \item \texttt{public void sumar(int a, int b)\{...\}}
\end{itemize}
\end{frame}

%-------------------------------------
\subsection{Comentarios en Java}

\begin{frame}[fragile]{3.3. Comentarios}
\textbf{Una línea}:
\begin{lstlisting}
// Comentario simple
\end{lstlisting}

\textbf{Multilínea}:
\begin{lstlisting}
/* Varias
   líneas
 */
\end{lstlisting}

\textbf{Documentación (javadoc)}:
\begin{lstlisting}
/**
 * Documento para describir la clase o método.
 * @param nombre
 * @return cadena de saludo
 */
\end{lstlisting}

\textbf{Ejercicio}:
Generar un pequeño programa con los tres tipos de comentarios y usar \texttt{javadoc}.
\end{frame}

%-------------------------------------
\subsection{Diferencias de sintaxis con C}

\begin{frame}{3.4. Diferencias principales}
\begin{itemize}
  \item \textbf{Sin funciones libres}: Todo método dentro de una clase.
  \item \textbf{Booleanos estrictos}: \texttt{true}, \texttt{false}, no enteros.
  \item \textbf{Sin aritmética de punteros}: Manejo de referencias.
  \item \textbf{Sin macros o preprocesador} (\texttt{\#define}, \texttt{\#include}).
  \item \textbf{Import} de paquetes (\texttt{java.util.*}, etc.) en vez de \texttt{\#include}.
\end{itemize}

\end{frame}

%=====================================
\section{Tipos de datos y Strings}

%-------------------------------------
\subsection{Tipos primitivos de datos}

\begin{frame}{4.1. Tipos primitivos}
\begin{itemize}
  \item \texttt{byte} (8 bits), \texttt{short} (16 bits), \texttt{int} (32 bits), \texttt{long} (64 bits)
  \item \texttt{float}, \texttt{double} (punto flotante de 32 y 64 bits)
  \item \texttt{char} (Unicode de 16 bits)
  \item \texttt{boolean} (\texttt{true}/\texttt{false})
\end{itemize}

\textbf{Ejemplo}:
\texttt{int contador = 0;}  
\texttt{double temperatura = 36.5;}

Siempre se pasan por \textbf{valor}.
\end{frame}

%-------------------------------------
\subsection{Variables y alcance}

\begin{frame}[fragile]{4.2.1. Declaración e inicialización}
\begin{lstlisting}[language=Java]
int contador = 0;
double altura = 1.75;
\end{lstlisting}

\textbf{4.2.2. Alcance (scope)}:
\begin{itemize}
  \item \textbf{Locales}: dentro de métodos.
  \item \textbf{Instancia}: definidas a nivel de clase, para cada objeto.
  \item \textbf{Clase (static)}: compartidas por todas las instancias.
\end{itemize}

\textbf{Ejemplo}:
\begin{lstlisting}[language=Java]
public class Alcance {
    int valor; // variable de instancia
    public void incrementar() {
        int incremento = 10; // variable local
        valor += incremento;
    }
}
\end{lstlisting}
\end{frame}

%-------------------------------------
\subsection{Operadores aritméticos, lógicos y relacionales}

\begin{frame}[fragile]{4.3. Operadores}
\begin{itemize}
    \item \textbf{Aritméticos}: `+ - * / %`  
    \item \textbf{Lógicos}: `&& || !`  
    \item \textbf{Relacionales}: `== != > < >= <=`
\end{itemize}

\textbf{Ejemplo}:
\begin{lstlisting}[language=Java]
int a = 7, b = 2;
System.out.println(a / b); // 3 (entero)
System.out.println(a % b); // 1
\end{lstlisting}

\textbf{Booleanos estrictos}:  
No se puede usar `if (a) {...}` como en C; debe ser `if (a > 0) {...}`.
\end{frame}

%-------------------------------------
\subsection{Clase String}

\begin{frame}[fragile]{4.4. Manejo de String}
\textbf{Creación y concatenación}:
\begin{lstlisting}[language=Java]
String saludo = "Hola";
String nombre = "Mundo";
String mensaje = saludo + ", " + nombre;
\end{lstlisting}

\textbf{Métodos comunes}:
\begin{itemize}
  \item \texttt{length()}, \texttt{charAt(i)}, \texttt{substring(i,j)}, \texttt{equals()}, ...
\end{itemize}

\textbf{Ejemplo}:
\begin{lstlisting}[language=Java]
System.out.println(mensaje.length());
System.out.println(mensaje.charAt(0));
\end{lstlisting}
\end{frame}

%-------------------------------------
\subsection{Inmutabilidad y comparación de String}

\begin{frame}[fragile]{4.5. Inmutabilidad}
\begin{itemize}
  \item Los objetos `String` no cambian su contenido.
  \item Al modificar, realmente se crea uno nuevo.
\end{itemize}
\begin{lstlisting}[language=Java]
String original = "Hola";
original = original + " Mundo"; 
// "Hola" sigue existiendo, "Hola Mundo" es otro objeto
\end{lstlisting}

\textbf{Comparación}:
\begin{lstlisting}[language=Java]
String a = "Hola";
String b = new String("Hola");
System.out.println(a == b);       // false (referencias distintas)
System.out.println(a.equals(b));  // true (contenido igual)
\end{lstlisting}
\end{frame}

%-------------------------------------
\subsection{StringBuilder y StringBuffer}

\begin{frame}[fragile]{4.6. Clases para manipular Strings (superavanzado)}
\textbf{StringBuilder}:
\begin{itemize}
  \item Mutable, rápido, no sincronizado.
\end{itemize}

\textbf{StringBuffer}:
\begin{itemize}
  \item Mutable, sincronizado, algo más lento.
\end{itemize}

\textbf{Ejemplo}:
\begin{lstlisting}[language=Java]
StringBuilder sb = new StringBuilder("Hola");
sb.append(" Mundo");
sb.insert(5, " querido");
System.out.println(sb.toString());
\end{lstlisting}
\end{frame}

%=====================================
\section{Manejo de memoria en Java vs. punteros en C}

%-------------------------------------
\subsection{Stack y Heap en Java}

\begin{frame}[fragile]{5.1. Modelo de memoria}
\textbf{Stack}:
\begin{itemize}
  \item Variables locales, referencias a objetos, llamadas de método.
\end{itemize}

\textbf{Heap}:
\begin{itemize}
  \item Objetos, variables de instancia.
  \item Liberados por el GC cuando no hay referencias.
\end{itemize}

\textbf{Ejemplo}:
\begin{lstlisting}[language=Java]
class Persona { String nombre; }
public class MemoriaDemo {
    public static void main(String[] args) {
        int edad = 25; // Stack
        Persona p = new Persona(); // Objeto en heap
        p.nombre = "Juan";
    }
}
\end{lstlisting}
\end{frame}

%-------------------------------------
\subsection{Referencias en Java vs punteros en C}

\begin{frame}{5.2. Referencias vs punteros}
\textbf{En Java}:
\begin{itemize}
  \item Referencias abstractas.
  \item Sin aritmética de punteros.
  \item Si es \texttt{null} y la usas, lanza \texttt{NullPointerException}.
\end{itemize}

\textbf{En C}:
\begin{itemize}
  \item Punteros directos a memoria.
  \item Pueden sumarse/restarse, apuntar a posiciones arbitrarias.
\end{itemize}

\textbf{Mayor seguridad en Java} a costa de menos control en bajo nivel.
\end{frame}

%-------------------------------------
\subsection{Garbage Collector}

\begin{frame}[fragile]{5.3. Recolección de basura}
\textbf{Funciona en segundo plano}:
\begin{itemize}
  \item Libera objetos sin referencias.
\end{itemize}

\textbf{Ventajas}:
\begin{itemize}
  \item Simplifica el código, evita fugas de memoria comunes en C.
\end{itemize}

\textbf{Limitaciones}:
\begin{itemize}
  \item No controlas exactamente cuándo se ejecuta el GC.
\end{itemize}

\textbf{Ejemplo}:
\begin{lstlisting}[language=Java]
String mensaje = new String("Hola");
mensaje = null; 
// "Hola" queda elegible para GC
\end{lstlisting}
\end{frame}

%-------------------------------------
\subsection{Paso de parámetros por valor}

\begin{frame}[fragile]{5.4. Paso por valor y referencias}
\textbf{Primitivos}: Copia del valor. 

\textbf{Objetos}: Copia de la referencia (apuntan al mismo objeto).

\begin{lstlisting}[language=Java]
public void cambiarNombre(Persona p) {
    p.nombre = "Ana";
}
Persona persona = new Persona();
persona.nombre = "Juan";
cambiarNombre(persona);
// persona.nombre -> "Ana"
\end{lstlisting}
\end{frame}

%-------------------------------------
\subsection{Comparativa con C}

\begin{frame}{5.5. Implicaciones de no tener punteros}
\textbf{Ventajas}:
\begin{itemize}
  \item Menos fugas de memoria, menos corrupción de punteros.
  \item Te centras en la lógica de negocio.
\end{itemize}

\textbf{Desventajas}:
\begin{itemize}
  \item Menos control sobre la gestión y optimización a bajo nivel.
  \item GC puede impactar en rendimiento en ciertos entornos.
\end{itemize}
\end{frame}

%=====================================
\section{Estructuras de control de flujo}

%-------------------------------------
\subsection{if, else if, else}

\begin{frame}[fragile]{6.1. Condicionales}
\begin{lstlisting}[language=Java]
if (condicion) {
   // ...
} else if (otraCond) {
   // ...
} else {
   // ...
}
\end{lstlisting}

\textbf{Igual a C}, pero en Java las condiciones deben ser booleanas (\texttt{true}/\texttt{false}).
\end{frame}

%-------------------------------------
\subsection{switch}

\begin{frame}[fragile]{6.2. switch vs C}
\textbf{Tipos permitidos}: \texttt{int}, \texttt{char}, \texttt{byte}, \texttt{short}, \texttt{String} (Java 7+).

\begin{lstlisting}[language=Java]
switch (dia) {
   case "Lunes":
      ...
      break;
   default:
      ...
}
\end{lstlisting}

\textbf{Nuevo switch en Java 12+} (expresión):
\begin{lstlisting}[language=Java]
String tipo = switch (dia) {
   case "Lunes","Martes" -> "Laboral";
   default -> "Fin de semana";
};
\end{lstlisting}
\end{frame}

%-------------------------------------
\subsection{Bucles for, while, do-while}

\begin{frame}[fragile]{6.3. Bucles}
\textbf{for}:
\begin{lstlisting}[language=Java]
for (int i=0; i<10; i++) {
   System.out.println(i);
}
\end{lstlisting}

\textbf{while}:
\begin{lstlisting}[language=Java]
int x=0;
while (x<10) {
   System.out.println(x);
   x++;
}
\end{lstlisting}

\textbf{do-while}:
\begin{lstlisting}[language=Java]
int y=0;
do {
   y++;
} while (y<10);
\end{lstlisting}
\end{frame}

%-------------------------------------
\subsection{break, continue, etiquetas}

\begin{frame}[fragile]{6.4. Instrucciones de interrupción}
\textbf{break}:
\begin{lstlisting}[language=Java]
for (int i=0; i<10; i++) {
  if (i==5) break; 
  ...
}
\end{lstlisting}

\textbf{continue}:
\begin{lstlisting}[language=Java]
for (int i=0; i<10; i++) {
  if (i % 2 == 0) continue;
  ...
}
\end{lstlisting}

\end{frame}

%=====================================
\section{Arrays y manipulación de datos}

%-------------------------------------
\subsection{Arrays unidimensionales}

\begin{frame}[fragile]{7.1. Declaración e inicialización}
\begin{lstlisting}[language=Java]
int[] numeros = new int[5];
int[] otros = {1,2,3,4,5};
System.out.println(numeros.length);
\end{lstlisting}

\textbf{En Java}:
\begin{itemize}
  \item Los arrays son objetos, con \texttt{length}.
  \item Tamaño fijo una vez creados.
  \item La memoria se reserva con \texttt{new}.
\end{itemize}
\end{frame}

%-------------------------------------
\subsection{Arrays multidimensionales}

\begin{frame}[fragile]{7.2. Bidimensionales y "jagged"}
\begin{lstlisting}[language=Java]
int[][] matriz = new int[3][4];
int[][] jagged = new int[3][];
jagged[0] = new int[4];
jagged[1] = new int[2];
jagged[2] = new int[5];
\end{lstlisting}

\textbf{En Java}: Realmente arrays de arrays.  

\textbf{En C}: \texttt{int mat[3][4];} (tamaño fijo).
\end{frame}

%-------------------------------------
\subsection{Recorrido de arrays}

\begin{frame}[fragile]{7.3. for y for-each}
\textbf{for normal}:
\begin{lstlisting}[language=Java]
for (int i=0; i<numeros.length; i++) {
    System.out.println(numeros[i]);
}
\end{lstlisting}

\textbf{for-each} (Java 5+):
\begin{lstlisting}[language=Java]
for (int num : numeros) {
    System.out.println(num);
}
\end{lstlisting}
\end{frame}

%-------------------------------------
\subsection{Arrays de objetos}

\begin{frame}[fragile]{7.4. Arrays de Strings u objetos}
\begin{lstlisting}[language=Java]
String[] nombres = {"Ana", "Luis", "Pedro"};
for (String n : nombres) {
    System.out.println(n);
}
\end{lstlisting}

En Java, \texttt{String} es un objeto.  
En C, serían punteros a \texttt{char}.
\end{frame}

%-------------------------------------
\subsection{Clase \texttt{java.util.Arrays}}

\begin{frame}[fragile]{7.5. Métodos útiles}
\begin{itemize}
  \item \texttt{Arrays.toString(numeros)}
  \item \texttt{Arrays.sort(numeros)}
  \item \texttt{Arrays.binarySearch(numeros, x)}
  \item \texttt{Arrays.fill(numeros, valor)}
\end{itemize}

\textbf{Ejemplo}:
\begin{lstlisting}[language=Java]
int[] nums = {5,3,8,1};
Arrays.sort(nums);
System.out.println(Arrays.toString(nums)); // [1,3,5,8]
\end{lstlisting}
\end{frame}

%=====================================
\section{Entrada/Salida (básica)}

%-------------------------------------
\subsection{Flujos de datos en consola}

\begin{frame}{8.1. Concepto de stream}
\textbf{System.out}: Salida estándar.  
\textbf{System.in}: Entrada estándar.  
\textbf{System.err}: Salida de errores.

Los flujos (streams) son secuencias de bytes. Para leerlos de forma sencilla en Java, usaremos \texttt{Scanner}.
\end{frame}

%-------------------------------------
\subsection{Uso de \texttt{System.out} y \texttt{System.err}}

\begin{frame}[fragile]{8.2. Salida en consola}
\begin{lstlisting}[language=Java]
System.out.println("Hola mundo!");
System.out.print("Sin salto de linea");
System.err.println("Mensaje de error");
\end{lstlisting}

\textbf{printf} para formatear (similar a C):
\begin{lstlisting}[language=Java]
int edad = 25;
System.out.printf("Tienes %d años\n", edad);
\end{lstlisting}
\end{frame}

%-------------------------------------
\subsection{Lectura desde consola}

\begin{frame}{8.3. \texttt{System.in}}
\texttt{System.in} es un InputStream “crudo”. Para leerlo cómodamente, usamos \textbf{Scanner}.
\end{frame}

%-------------------------------------
\subsection{Clase Scanner}

\begin{frame}[fragile]{8.4. Ejemplo con Scanner}
\begin{lstlisting}[language=Java]
import java.util.Scanner;

Scanner entrada = new Scanner(System.in);
System.out.println("Introduce tu nombre: ");
String nombre = entrada.nextLine();

System.out.println("Introduce tu edad: ");
int edad = entrada.nextInt();

System.out.println("Hola " + nombre + 
                   ", tienes " + edad + " años.");
entrada.close();
\end{lstlisting}
\end{frame}

%-------------------------------------
\subsection{Formateo de salida}

\begin{frame}[fragile]{8.5. \texttt{printf} y \texttt{format}}
\textbf{Ejemplo}:
\begin{lstlisting}[language=Java]
double precio = 12.5;
System.out.printf("El precio es %.2f euros\n", precio);
System.out.format("El precio es %.2f euros\n", precio);
\end{lstlisting}

\textbf{Ejercicio}:  
Leer nombre y edad, y mostrar:  
\texttt{[Usuario: Juan, Edad: 30 años]}
\end{frame}

\end{document}
